\documentclass[aspectratio=169,t,11pt,table]{beamer}
\usepackage{slides,math}

%\input{preambles/header_beamer}

\graphicspath{{../}}

% Optionally define `accent`/`accent2` colors for theme customization
% I recommend changing the top slider on this: https://hslpicker.com/#1e9400
\definecolor{accent}{HTML}{940034}
\definecolor{accent2}{HTML}{006896}

\title{Generative AI at Work}
% originally this was created for the MIT Future of Work Conference in December of 2024
% was presented at Deutsche Bank and at the TASK conference 
\date{September 2025}
\author{Lindsey Raymond }
\institute{ { Microsoft Research \& MIT}  \\
TASKS VII: The Economic Impacts of AI on Work and Labor Markets }
% \\ {\small  Joint with} \\  {\small Erik Brynjolfsson (Stanford) and Danielle Li (MIT)}
%\addbibresource{references.bib}
\addbibresource{AugI_Bibliography.bib}

\begin{document}

% ------------------------------------------------------------------------------
\begin{frame}[noframenumbering,plain]
\maketitle

%\bottomleft{\footnotesize $^*$A bit of extra info here. Add an asterich to title or author}
\end{frame}
% ------------------------------------------------------------------------------

% ------------------------------------------------------------------------------
% turn this on to turn on an automatic outline slide \section{Common Items}
% ------------------------------------------------------------------------------


\begin{frame}{The old: programmable computers}
\medskip
\bi  
\item[] Traditionally, computers require inputs ($X$) and \underline{instructions} ($f(X)$) to produce output ($Y$)
\ei 
\begin{figure}
\begin{center}
    \includegraphics[scale=0.60]{figures/presentation_images/trad_programming_LR.png}
\end{center}
\end{figure}
\bi  
\item[] Computers limited to tasks where we know $f(X)$
\ei 
\end{frame}


\begin{frame}{The new: learn f(X) from data}
\medskip
\bi  
\item[] Machine learning requires inputs ($X$) and output ($Y$) to \underline{learn} instructions ($f(X)$)
\ei 
\begin{figure}
\begin{center}
    \includegraphics[scale=0.60]{figures/presentation_images/ml_programming_LR.png}
\end{center}
\end{figure}
\bi  
\item $f(X)$ is an output, not an input
\item Polyani's paradox (\cite{autorPolanyiParadoxShape2014})
\ei 
\end{frame}

% add pros/cons around computers playing a role in tasks 
\begin{comment}
\begin{frame}{Generative AI: optimism and apprehension}
\alert{Generative AI:} ML that uses learned patterns to generate new content
\smallskip
\begin{enumerate}
    \item Can generative AI expand the set of tasks machines can perform?
    \bi
    \medskip
    \item Reasons to be skeptical: 
    \bi 
    \item Decision support tools + human seem to do worse than human or algorithm alone
    \medskip
    \item Gulf between lab and real world
    \medskip
    \item Coherent but inaccurate 
    \ei 
    \ei  \pause
    % heterogeneity in effects -> f(X) unevenly distributed 
    \bigskip
    \item Will generative AI also change the experience of work?

\end{enumerate}
\end{frame}
\end{comment}

\begin{frame}{Generative AI: optimism and apprehension}
\alert{Generative AI:} ML that uses learned patterns to generate new content
\smallskip
\begin{itemize}
    \item \textbf{Optimism}
    \begin{itemize}
        \item Expands the set of tasks machines can perform
        \medskip
    \end{itemize}
    \medskip
    \item \textbf{Apprehension}
        \begin{itemize}
            \item Adjustment costs, gulf between lab and real-world scenarios
        \medskip
        \pause 
        \item Humans + machines 
        \begin{itemize}
            \item Decision support tools + human seem to do worse than human or algorithm alone (Vaccaro, Almaatouq \& Malone, 2024)
        \end{itemize}
        \medskip
    \item Coherent but inaccurate
    \end{itemize}
\end{itemize}
\end{frame}

\begin{frame}{A paper: Generative AI at Work} {Joint work with Erik Brynjolfsson (Stanford) and Danielle Li (MIT)}
\begin{enumerate}
    \item Does access to generative AI suggestions increase productivity? 
    
    \only<3>{ \alert{15\% increase in resolutions per hour} }
      \medskip
    \item For which workers and why? 

    \only<3>{\alert{Increases are driven by less experienced, less skilled workers} }
    \medskip
    \item How does this change the experience of work? 
        \only<3>{\alert{Attrition, customer interactions}}
    \medskip
\end{enumerate}
\medskip 
\only<2>{\textcolor{teal}{Setting}: Customer support  
	\bigskip
 
\textcolor{teal}{Technology}: AI assistant built on GPT-3  \pause
	\bigskip
 
\textcolor{teal}{Empirical design}: Staggered roll-out  at a Fortune 500  software company
\medskip }
\end{frame}



\begin{frame}{Agenda}
\be
\item \textbf{Setting: customer support}
\bi 
\item AI chatbot design and output
\ei 
\medskip
\item Data and study design
\medskip 
\item Productivity effects
\bi
\item Aggregate and by skill and experience
\item Productivity-experience curve
\item Copying or learning?
\item Learning what?
\ei 
\medskip
\item \textcolor{gray}{Experience of work}
\bi 
\item \textcolor{gray}{Customer interactions}
\item \textcolor{gray}{Attrition}
\ei 
\ee 
\end{frame}

\begin{frame}{Customer support is leading AI adoption}
\textcolor{raspberry}{Technical feasibility: }
\begin{itemize}
\item Millions of chats with automatically labeled outcomes 
\end{itemize} 
\medskip \pause
\textcolor{raspberry}{Business need}:
\begin{itemize} 
\item High turnover %(60 percent)
\item Large and persistent performance differences  
\end{itemize} 
\medskip \pause
\textcolor{raspberry}{Technical and social skills required: }
\begin{itemize}
\item Computer and accounting knowledge
\item Social and emotional intelligence 
\end{itemize} 
\end{frame}

\begin{frame}{Customer support}

\begin{figure}
\centering
\makebox[\linewidth]{\includegraphics[scale=0.23]{figures/presentation_images/live-engagescreen1.pdf}}
\end{figure}
\end{frame}

%%%%%DESCRIBE WHY CUSTOMER SERVICE IS AN ML PROBLEM
\begin{frame}{Why is customer support an ML problem?}
\begin{enumerate}
\item \alert{Diagnose the problem} \smallskip
    \begin{itemize}
    \item What underlying problems present as ``I can't login"? \bigskip
     \end{itemize}
    \pause 
 \item \alert{Predict what will resolve it} \smallskip
\begin{itemize}
 \item What types of solutions have typically worked for this problem? \bigskip
  \end{itemize} 
 \pause 
  \item \alert{Recognize what gets people yelled at} \smallskip
   \begin{itemize}
 \item How to communicate in ways that alleviate customer frustration?
  \end{itemize}
\end{enumerate}

\end{frame}


%%%%%DESCRIBE WHY CUSTOMER SERVICE IS AN ML PROBLEM
\begin{frame}{Our AI system}
\textcolor{raspberry}{Model design}: 
\begin{enumerate}
\item \alert{Start} from a large language model from the Generative Pre-Trained Transformer (GPT) family of models (GPT-3)
\medskip \pause
\item \alert{Train} model to focus on customer service applications using data on historical chats  
\bi
\item Text of conversation
\item Labelled with outcomes: call duration, successfully resolved, customer satisfaction, ``high performer'' or not
\ei 
\pause
 \item \alert{Fine-tuned} with human feedback
\end{enumerate}
\bigskip \pause
\textcolor{raspberry}{System function}: 
\bi
\item Real-time text suggestions and links to technical material
\item Worker retains \alert{full discretion}
\ei
\end{frame}




\begin{comment}
%%%%%DESCRIBE OUR SETTING
\begin{frame}{Our tool}
% https://latexdraw.com/how-to-annotate-an-image-in-latex/
\begin{figure}
\begin{center}
    \includegraphics[scale=0.40]{figures/Setting/cresta_customer_v2_redacted.png}
\end{center}
\end{figure}
\end{frame}
\end{comment}

\begin{frame}{AI tool provides suggestions to the agent}
% https://latexdraw.com/how-to-annotate-an-image-in-latex/
\begin{figure}
\begin{center}
\makebox[\linewidth]{
\begin{tabular}{cc}
\includegraphics[scale=0.25]{figures/productimages/cresta_customer_v2_redacted.png} & 
 \includegraphics[scale=0.25]{figures/productimages/cresta_example_suggestion.png}
\end{tabular}
}
\end{center}
\end{figure}
\begin{itemize}
\item Recommendations based on responses that are most correlated with successful outcomes
\item In this case: establishing a friendly, reassuring rapport.
\end{itemize} 
\end{frame} 




\begin{comment}
%\section{Data and Study Overview}
% hyperlink to the deployment graph
%\frame{\tableofcontents[currentsubsection]}
\begin{frame}[label=study-design]{Data and study design}
\medskip
\alert{Data}
\begin{itemize}
%\item \textcolor{blue}{Research Question:}  \smallskip
\item \textbf{Conversations:} 3,000,000 chats between January 2020 and June 2021 
\item \textbf{Agents:} 3,000 agents, 140 teams and 5 firms 
\item \textcolor{raspberry}{Data:} Chat text, AI output, agent interactions 
\item \textcolor{raspberry}{Outcomes:} Issues solved per hour and customer satisfaction
\end{itemize} 

\end{frame}
 \begin{frame}{Data and study design}
\alert{Study Design}
\begin{equation*} \label{eq:dd_main}
y_{it}= \delta_{t} + \alpha_{i} + \sum_{k \ne -1} \beta_{k}\times 1[t=k] \times AI_{it} + \gamma X_{it}+ \epsilon_{it}
\end{equation*}
\begin{itemize} 
\item AI rolled out over six months at the agent level \hyperlink{deployment-timeline}{\beamerbutton{Timeline}}  
\item $AI_{it}$ indicates agent $i$ has access to AI assistance at time $t$
\item Use estimators robust to differential timing
\end{itemize}

\end{frame}   
\end{comment}

\begin{frame}{Agenda}
\be
\item \textcolor{asher}{Setting: customer support}
\bi 
\item \textcolor{asher}{AI chatbot design and output}
\ei 
\medskip
\item \textbf{Data and study design}
\medskip 
\item Productivity effects
\bi
\item Aggregate and by skill and experience
\item Productivity-experience curve
\item Copying or learning?
\item Learning what?
\ei 
\medskip
\item Experience of work
\bi 
\item Customer interactions
\ei 
\ee 
\end{frame} 


\begin{frame}[label=data]{Our data}
\textbf{Sample}
\bi
\footnotesize
\item 3 million customer service chats between 2020-2021 
\item 5,000 agents; 1,500 treated
\item 5 firms (data firm and 4 outsourcing firms)
\ei
\medskip
\textbf{Chat level information}
\bi
\footnotesize
\item Customer and agent text
\item Adherence to AI suggestions
\item Chat duration
\ei 
\medskip
\textbf{Agent-month level information}
\bi
\footnotesize
\item Average call resolution rate 
\item Average surveyed customer satisfaction
\ei
\hyperlink{summary-stats}{\beamerbutton{Summary stats}} 
\hyperlink{balance}{\beamerbutton{Balance table}} 
\end{frame}


\begin{frame}{Empirical strategy: staggered rollout of AI}
\begin{columns}[T] % top align the columns 
    \begin{column}{.4\textwidth}
\textcolor{digitalblue}{Rollout shaped by two constraints:}
\begin{itemize}
    \item Scheduling
    \begin{itemize}
        \item Limited training personnel $\implies$ long rollout
        \item Scheduling constraints $\implies$ agent-level rollout
    \end{itemize}
    \vspace{2mm}
    \item Budget
        \begin{itemize}
        \item Small RCT before deployment
        \item Fixed dollar budget
    \end{itemize}
\end{itemize}
\end{column}
    \hspace{-4mm}
     \begin{column}{.6\textwidth}
     \begin{center}
\includegraphics[scale=0.55]{figures/24June2024/deployment_timeline_cumulative.pdf}
\end{center}
     \end{column}
        \hfill
  \end{columns}
\end{frame}


 \begin{frame}[label=studydesign]{Empirical strategy: staggered rollout of AI}
\alert{Study Design}
\only<1>{ 
$$ \textcolor{raspberry}{\underbrace{\bm{y_{it}}}_{\makebox[0pt]{\text{\shortstack{\normalsize Resolutions per hour }}}}} = \alpha_{i} + \delta_{t} + 
     \sum^{J}_{j\ne-1}\beta_j (AI_{it} \times  \Ind[t=j]) + \gamma X_{it} + \epsilon_{it}  $$ 
}
\only<2>{ 
$$ y_{it} = \textcolor{raspberry}{\underbrace{\bm{\alpha_{i} + \delta_{t}}}_{\makebox[0pt]{\text{\shortstack{\normalsize Agent + \\ Month FE }}}}} + 
     \sum^{J}_{j\ne-1}\beta_j (AI_{it} \times  \Ind[t=j]) + \gamma X_{it} + \epsilon_{it}  $$ 
}
\only<3>{ 
$$ y_{it} = \alpha_{i} + \delta_{t} + 
     \sum^{J}_{j\ne-1}\beta_j (\textcolor{raspberry}{\underbrace{\bm{AI_{it}}}_{\makebox[0pt]{\text{\shortstack{\normalsize AI \\ Turned On }}}}} \times  \Ind[t=j]) + \gamma X_{it} + \epsilon_{it}  $$ 
}
\only<4->{ 
$$ y_{it} = \alpha_{i} + \delta_{t} + 
     \sum^{J}_{j\ne-1}\beta_j (AI_{it} \times  \Ind[t=j]) + \gamma \textcolor{raspberry}{\underbrace{\bm{X_{it}}}_{\makebox[0pt]{\text{\shortstack{\normalsize Controls }}}}} + \epsilon_{it}  $$ 
}
\begin{itemize} 
\item AI rolled out over six months at the agent level  
\item $AI_{it}$ indicates agent $i$ has access to AI assistance at time $t$
\item $ X_{it}$ months of agent tenure
\item Use estimators robust to differential timing
\item Data available at the agent-month level
\end{itemize}

\end{frame} 




% replace agenda 
\begin{frame}{Agenda}
\be
\item \textcolor{asher}{Setting: customer support}
\bi 
\item \textcolor{asher}{AI chatbot design and output}
\ei 
\medskip
\item \textcolor{asher}{Data and study design}
\medskip 
\item \textbf{Productivity effects}
\bi
\item Aggregate and by skill and experience
\item Productivity-experience curve
\item Copying or learning?
\item Learning what?
\ei 
\medskip
\item \textcolor{gray}{Experience of work}
\bi 
\item \textcolor{gray}{Customer interactions}
\ei 
\ee 
\end{frame} 

\begin{comment}
    

\begin{frame}{Gains are evident in the raw data}
\begin{figure}
\centering
\makebox[\linewidth]{
\begin{tabular}{cc}
\includegraphics[height=3.6cm]{figures/24June2024/productivity_rph_cresta_kdensity_0.pdf} & \pause \includegraphics[height=3.6cm]{figures/24June2024/productivity_aht_kdensity_0.pdf} \\

\pause \includegraphics[height=3.6cm]{figures/24June2024/productivity_cph_kdensity_0.pdf} &  \pause \includegraphics[height=3.6cm]{figures/24June2024/productivity_fcr_cresta_kdensity_0.pdf} \\
\end{tabular}
}
\end{figure}
\end{frame}
\end{comment}

\begin{frame}{Gains are evident in the raw data}
\begin{figure}
\centering
\makebox[\linewidth]{
\begin{tabular}{cc}
\includegraphics[height=3.6cm]{figures/24June2024/productivity_rph_cresta_kdensity_1.pdf} &  \pause \includegraphics[height=3.6cm]{figures/24June2024/productivity_aht_kdensity_1.pdf} \\ \pause

\includegraphics[height=3.6cm]{figures/24June2024/productivity_cph_kdensity_1.pdf} &  \pause \includegraphics[height=3.6cm]{figures/24June2024/productivity_fcr_cresta_kdensity_1.pdf} \\
\end{tabular}
}
\end{figure}
\end{frame}



\begin{frame}[label=dd-figure]{On average, a 15 percent productivity improvement}
\begin{figure}[ht!]
\begin{center}
\makebox[\linewidth]{
  \includegraphics[scale=0.5]{figures/15Aug2024/ES_rph_cresta_AS.pdf} 
}
\end{center}
\end{figure}   
\hyperlink{dd-table}{\beamerbutton{Difference-in-differences table}} 
\hyperlink{adherence}{\beamerbutton{AI adherence}} 
\end{frame}


\begin{frame}{Similar results across estimators}
\begin{figure}[ht!]
\begin{center}
\makebox[\linewidth]{\includegraphics[scale=0.7]{figures/20Jul2023/ALL_ES_rph_cresta.pdf}
}
\end{center}
\end{figure}
\end{frame}


\begin{frame}[label=adherence]{35\% of AI suggestions followed}
\begin{figure}[ht!]
\begin{center}
\makebox[\linewidth]{
  \includegraphics[scale=0.7]{figures/24June2024/receptivity_dist.pdf}
}
\end{center}
\end{figure}    
\end{frame}


\begin{frame}{Event studies, additional outcomes}
\begin{figure}[ht!]
\begin{center}
\makebox[\linewidth]{
\begin{tabular}{cc}
\textsc{A. Average Handle Time} & \textsc{B.  Chats Per Hour}\\
\includegraphics[scale=0.25]{figures/15Aug2024/ES_aht_AS.pdf} & \pause \includegraphics[scale=0.25]{figures/15Aug2024/ES_cph_AS.pdf} 
\\
\textsc{C. Resolution Rate} & \textsc{D. Customer Satisfaction (NPS)}\\
\pause  \includegraphics[scale=0.25]{figures/15Aug2024/ES_fcr_cresta_AS.pdf} & \pause  \includegraphics[scale=0.25]{figures/15Aug2024/ES_tnps_cresta_AS.pdf} 
\\
\end{tabular}
}
\end{center}
\end{figure}
\end{frame}




%%%%%%%%%%%%%DIFFERENTIAL EFFECTS BY WORKER SKILL LEVEL
\begin{frame}{How does the value of AI assistance vary across workers?}

\alert{Higher ability workers use AI more effectively}
\bi
\item Prior evidence of skill bias in IT \citep{bartel2007}
\smallskip
\item AI requires discretion 
\ei
\bigskip \pause
\alert{Lower-ability workers have more knowledge gaps}
\bi
\item ML learns from best workers
\smallskip
\item Better workers already know AI suggestions, less skilled workers do not
\ei
\bigskip \pause
\alert{We measure worker skill and experience}
\bi
\item \textcolor{teal}{Skill}: pre-AI index of all outcomes
\smallskip
\item \textcolor{teal}{Experience}: pre-AI tenure with the firm
\ei
\end{frame}

\begin{frame}{Highest returns for the lowest skill agents}

\begin{figure}
\centering
\makebox[\linewidth]{
\begin{tabular}{c}
\includegraphics[scale=0.55]{figures/15Aug2024/DD_rph_cresta_skill_tenure.pdf} 
\end{tabular}
}
\end{figure}
\hyperlink{tenure1}{\beamergotobutton{Tenure}}
\end{frame}



\begin{frame}[label=skill]{Additional outcomes by worker skill}
\begin{figure}
\begin{center}
\makebox[\linewidth]{
\begin{tabular}{cc}
\textsc{A. Average Handle Time} & \textsc{B.  Chats Per Hour}\\
\includegraphics[scale=0.26]{figures/15Aug2024/DD_aht_skill_tenure.pdf} \pause & \includegraphics[scale=0.26]{figures/15Aug2024/DD_cph_skill_tenure.pdf}  \pause
\\
\textsc{C. Resolution Rate} & \textsc{D.  Customer Satisfaction (NPS)}\\
\includegraphics[scale=0.26]{figures/15Aug2024/DD_fcr_cresta_skill_tenure.pdf} \pause & \includegraphics[scale=0.26]{figures/15Aug2024/DD_tnps_cresta_skill_tenure.pdf} 
\\
\end{tabular}
}
\end{center}
\end{figure}
\end{frame}

%%%%%%%%%%%%%DIFFERENTIAL EFFECTS BY WORKER TENURE 

\begin{frame}[label=tenure1]{Highest returns for the newest agents}
\begin{figure}
\centering
\makebox[\linewidth]{
\begin{tabular}{c}
\includegraphics[scale=0.6]{figures/15Aug2024/DD_rph_cresta_tenure_skill.pdf} 
\end{tabular}
}
\end{figure}
%\hyperlink{skill}{\beamerbutton{Back}} 
\end{frame}


\begin{frame}[label=tenure2]{Additional outcomes by worker tenure}
\begin{figure}
\begin{center}
\makebox[\linewidth]{
\begin{tabular}{cc}
\textsc{A. Average Handle Time} & \textsc{B.  Chats Per Hour}\\
\includegraphics[scale=0.26]{figures/15Aug2024/DD_aht_tenure_skill.pdf} & \pause \includegraphics[scale=0.26]{figures/15Aug2024/DD_cph_tenure_skill.pdf}  \pause
\\
\textsc{C. Resolution Rate} & \textsc{D.  Customer Satisfaction (NPS)}\\
\includegraphics[scale=0.26]{figures/15Aug2024/DD_fcr_cresta_tenure_skill.pdf}  \pause & \includegraphics[scale=0.26]{figures/15Aug2024/DD_tnps_cresta_tenure_skill.pdf} 
\\
\end{tabular}
}
\end{center}
\end{figure}
%\hyperlink{skill}{\beamerbutton{Back}} 
\end{frame}


\begin{frame}{AI assistance helps newer agents ``catch up''}
\begin{figure}
\begin{center}
\makebox[\linewidth]{
\begin{tabular}{c}
  \begin{overprint}
 \onslide<1>\includegraphics[scale=0.6]{figures/15Aug2024/learningcurve_rph_cresta_0.pdf}
 \onslide<2>\includegraphics[scale=0.6]{figures/15Aug2024/learningcurve_rph_cresta_1.pdf}
 \onslide<3>\includegraphics[scale=0.6]{figures/15Aug2024/learningcurve_rph_cresta.pdf}
    \end{overprint} 
\end{tabular}
}
\end{center}
\end{figure}
\smallskip
\end{frame}



% add subbuttons for these 
\begin{frame}[label=robustness]{Robustness checks}
\begin{itemize}
    \item Instrument individual adoption with earliest date of team adoption \hyperlink{teamcompanyiv}{\beamergotobutton{IV estimates}} 
    \vspace{3mm}
    \item Alternative diff-in-diff specifications \hyperlink{altspecs}{\beamergotobutton{Alternative specifications}}
    \vspace{3mm}
    \item Analyze randomized control trial \hyperlink{rct}{\beamergotobutton{Randomized control trial}}
    \vspace{3mm}
    \item Vary the level of clustering \hyperlink{clustering}{\beamergotobutton{Clustering}}
    \vspace{3mm}
    \item Weight by the number of chats \hyperlink{weighting}{\beamergotobutton{Weighting}}
    \vspace{3mm}
        \item Mean reversion \hyperlink{meanrev}{\beamergotobutton{mean reversion?}}
\end{itemize}
\end{frame}

%%%%%%%%%%%%%%%%%%%%%%
%%%%%%%%%%%%%%%%%%%%%%
% results we can give ourself some credit for! 
% add a slide lead up to 
% pass thtough the text 

\begin{frame}{Testing for worker learning}
\begin{enumerate}
    \item \alert{Learning:} examples help workers learn how to manage customers and how to solve problems
  \pause 
  \bigskip
    \begin{itemize} 
    \item Manager coaching is constrained by time and human capabilities
    \item[] $\implies$ AI-exposed workers will do better even without AI access
    \end{itemize} 
    \pause  
    \bigskip
    \item \alert{Blind copying:} workers blindly follow AI recommendations and develop no more human capital than otherwise
     \bigskip
    \begin{itemize} \pause 
    \item[] $\implies$ AI-exposed workers will be no better without AI access
    \end{itemize} 
\end{enumerate}

\end{frame}


\begin{frame}{Occasional periods of AI outages}
\begin{figure}
\begin{center}
\makebox[\linewidth]{
\begin{tabular}{c}
  \includegraphics[scale=0.6]{figures/24June2024/outage3_10sept2020_graph.pdf}
\end{tabular}
}
\end{center}
\end{figure}
\smallskip
\end{frame}

\begin{frame}{AI-exposed workers faster during AI outages after exposure}
\begin{figure}
\centering
\makebox[\linewidth]{
\begin{tabular}{cc}
a.  AI Working & b.  AI Outage \\
\includegraphics[scale=.55]{figures/15Aug2024/ES_aht_nonoutage.pdf} \pause  & \includegraphics[scale=.55]{figures/15Aug2024/ES_aht_outage.pdf} \\
\end{tabular}
}
\end{figure}
\end{frame}


\begin{frame}{Intensive AI users are faster during outages}
\begin{figure}
\centering
\makebox[\linewidth]{
\begin{tabular}{cc}
a.  High Initial Adherence & b.  Low Initial Adherence \\
\includegraphics[scale=.55]{figures/15Aug2024/ES_aht_outage_recept3.pdf} \pause  & \includegraphics[scale=.55]{figures/15Aug2024/ES_aht_outage_recept1.pdf} \\
\end{tabular}
}
\end{figure}
%\begin{itemize}
%    \item Consistent with other work from the education context
%\end{itemize}
\end{frame}


\begin{frame}{Assessing mechanism}
\begin{enumerate}
    \item AI learns patterns associated with success
        \begin{itemize} 
    \item[] \alert{Challenge:} patterns are not human-interpretable
    \end{itemize} 
    \medskip \pause
    \item Chat transcripts are a window into what workers do
    \begin{itemize} 
    \item[] $\implies$ Assess chat text
    \end{itemize} 
    \medskip  \pause
      \end{enumerate}
    \textcolor{raspberry}{Evidence on mechanism:}
    \begin{itemize}
    \item Analyze impacts on English proficiency
    \item \textcolor{gray}{Compare high and low skill workers' language convergence}
    \item Impacts by problem topic
      \end{itemize}
\end{frame}



\begin{frame}{English language fluency}
\begin{itemize}
    \item Good call center agents communicate with \alert{comprehensible} and \alert{native-sounding} English
    \vspace{2mm}
    \item We use the chat transcripts to measure;
    \begin{enumerate}
        \item \textbf{Comprehensibility} - a score ranging from 1 (``very difficult to comprehend'') to 5 (``very fluent and easily understandable, with no significant errors'')
        \vspace{2mm}
        \item\textbf{Native fluency} -  Interagency Language Roundtable ``functionally native'' proficiency standard score from 1 to 5 
    \end{enumerate}
    \item Transcripts scored by Gemini, a LLM, and LLM scores are human-validated
    \vspace{2mm}
    \item Does access to AI impact agent comprehensibility or native fluency?
\end{itemize}
\end{frame}

\begin{frame}{AI access increases English fluency}
\begin{figure}
\begin{center}
\makebox[\linewidth]{
    \begin{tabular}{cc}
\textsc{Comprehensibility Score} &  \textsc{Native Fluency Score}\\
    \includegraphics[scale=0.4]{figures/15Aug2024/ES_fluency_score_AS.pdf} & \pause
  \includegraphics[scale=0.4]{figures/15Aug2024/ES_native_score_AS.pdf}\\ 
\end{tabular}
}
\end{center}
\end{figure}
\begin{itemize}
    \item 5\% increase in comprehensibility score
    \item 6\% increase in native fluency score
\end{itemize}
\end{frame}

\begin{frame}{Filipino agents linguistic improvements double the U.S workers}
\begin{figure}
\begin{center}
\makebox[\linewidth]{
    \begin{tabular}{cc}
\textsc{Comprehensibility Score} & \textsc{Native Fluency Score}\\
\includegraphics[scale=0.33]{figures/15Aug2024/languagedd_bycountry_comp.pdf} & \pause \includegraphics[scale=0.33]{figures/15Aug2024/languagedd_bycountry_native.pdf} 
\end{tabular}
}
\end{center}
\end{figure}
\begin{itemize}
\item \textbf{Comprehensibility:} US +3.5\%; Philippines +5.7\%  
\item \textbf{Native Fluency:} US +2.9\%; Philippines +5.3\%
\item Customer engagement occurs across a wide range of jobs
\end{itemize}
\end{frame}



\begin{frame}{Identify the chat topic}
\begin{columns}[T] % top align the columns 
\begin{column}{.4\textwidth}
\begin{center}
    \begin{itemize}
        \item Randomly sample 5,000 chats
        \item Ask for a short description of chat topic
        \item Cluster descriptions into 50 topics
        \item Classify all conversations into topic 
        \item Human validation
    \end{itemize}
\end{center}
     \end{column}
      \hspace{-4mm}
    \begin{column}{.6\textwidth}
 \begin{center}
\includegraphics[scale=0.21, angle=90]{figures/24June2024/hist_topic.pdf}
\end{center}
\end{column}
        \hfill
  \end{columns}
\end{frame}



\begin{comment}
  \begin{frame}[label=topicdesign]{Empirical strategy for chat topic analysis}
\alert{Study Design}
$$ y_{itc} = \delta_{t} + \alpha_{i} + \sum_{r=1}^{4} \beta_r (AI_{it} \times Topic_{cr}) + \gamma X_{itc} + \epsilon_{itc} $$ 

\begin{itemize}
    \item $\beta_r$ impact of AI assistance on call in duration for chats in topic frequency category $r$
    \item Data is at the chat level
    \item $AI_{it}$ is an indicator for if agent $i$ has access to AI recommendations in year-month $t$
    \item $X_{it}$ includes fixed effects for agent tenure and the overall category of topic frequency 
\end{itemize}
\end{frame}    
\end{comment}



\begin{frame}{Impacts of AI along topic distribution}
% \only<1>{
%\begin{enumerate}
 %   \item Quality of AI predictions improves with data, but humans also have experience
%\end{enumerate} }
%\vspace{2mm}
     \centerline{\includegraphics[scale=0.65]{figures/15Aug2024/chat_duration_minutes_byagg_topic_freq_pct.pdf}}
   % \only<3->{\centerline{\includegraphics[scale=0.65]{figures/15Aug2024/chat_duration_byindivtopicexp_pct.pdf}}}
\begin{center}
\end{center}
\begin{itemize}
\item Effects of AI on novel problems
\item Pulling workers towards the middle of the performance distribution?
\end{itemize}
\end{frame} 

\begin{comment}
\begin{frame}{Assessing what workers are learning}
\begin{itemize}
    \item Top performers generate successful outcomes
    \vspace{2mm}
        \begin{itemize} 
    \item[] $\implies$ Compare high and low skill workers conversations before and after AI 
    \end{itemize} 
    \vspace{2mm}
    \item Do we see changes in workers' language around the introduction of AI assistance?
    \vspace{2mm}
    \item Does high and low skill workers' language converge?
    \vspace{2mm}
    \end{itemize}
\end{frame}


    

\begin{frame}{AI changes worker language}
\begin{figure}
\begin{center}
\makebox[\linewidth]{
\begin{tabular}{c}
\includegraphics[scale=0.6]{figures/06Oct2023/within_agent_sim_to_month_minus1.pdf}
\end{tabular}
}
\end{center}
\end{figure}
\end{frame}


\begin{frame}[label=high-low-gap]{After AI access, low-skill workers sound more like high-skill workers}
\begin{figure}
\begin{center}
\makebox[\linewidth]{
\begin{tabular}{c}
\includegraphics[scale=0.75]{figures/05Apr2023/HL_chat_sim_nagents.pdf}
\end{tabular}
}
\end{center}
\end{figure}
\end{frame}
\end{comment}

\begin{frame}{Conclusion}
\alert{Evidence from contact centers}
\begin{itemize}
\item Clear, quantifiable outcomes  
\item Mix of repetitive tasks and less routine tasks
\end{itemize}

\alert{Raises further questions:}
\begin{itemize}
\item \textbf{Employment:} Corresponds to 12\% fewer worker-hours
\begin{itemize}
    \item Where will those jobs be located?
\end{itemize}
\item \textbf{Wages:} Nature of the existing contract structure
\begin{itemize}
    \item Workers: performance-based pay
    \item Outsourcing firms: contracts in hours of work
\end{itemize}
\end{itemize}

\bigskip
\textcolor{raspberry}{Questions:} lraymond@mit.edu
\end{frame}


% ------------------------------------------------------------------------------
\appendix
\begin{frame}[allowframebreaks,noframenumbering]{Appendix}
  \thispagestyle{empty}
  %\printbibliography
\end{frame}


\begin{frame}{Sample Summary Statistics}\label{summary-stats}
\begin{table}
\scalebox{.70}{\makebox[\linewidth]{\input{tables/12Mar2023/summary_table}}}
    \caption{Sample Summary Statistics}
    \label{tab:summarystats}
\end{table}
\end{frame}
% ------------------------------------------------------------------------------

\begin{frame}[label=dd-table]{On average, a 15 percent productivity improvement}
\begin{figure}[ht!]
\begin{center}
\scalebox{.6}{\makebox[\linewidth]{\input{tables/15Aug2024/maindd_rph.tex}}}
\end{center}
\end{figure}   
\hyperlink{dd-figure}{\beamerbutton{Back}}
\end{frame}

\begin{frame}[label=teamcompanyiv]{Team Company IV}
\begin{center}
\scalebox{.6}{\makebox[\linewidth]{\input{tables/15Aug2024/ivdd_companylocteam.tex}}}
\end{center}
\hyperlink{robustness}{\beamerbutton{Back}} 
\end{frame}

\begin{frame}[label=altspecs]{Alternative Specifications}
\begin{center}
\scalebox{.6}{\makebox[\linewidth]{\input{tables/15Aug2024/refrequest_maindd_rph.tex}}}
\end{center}
\hyperlink{robustness}{\beamerbutton{Back}} 
\end{frame}





\begin{frame}[label=rct]{Randomized control trial}
\begin{center}
\scalebox{.75}{\makebox[\linewidth]{\input{tables/15Aug2024/rctdd.tex}}}
\end{center}
\hyperlink{robustness}{\beamerbutton{Back}} 
\end{frame}

\begin{frame}[label=clustering]{Robustness to alternative clustering}
\begin{center}
\scalebox{.75}{\makebox[\linewidth]{\input{tables/15Aug2024/robustness_clustering.tex}}}
\end{center}
\hyperlink{robustness}{\beamerbutton{Back}} 
\end{frame}


\begin{frame}[label=weighting]{Weighting by number of chats}
\begin{center}
\scalebox{.75}{\makebox[\linewidth]{\input{tables/15Aug2024/robustdd_nchat_weight.tex}}}
\end{center}
\hyperlink{robustness}{\beamerbutton{Back}} 
\end{frame}

\begin{frame}[label=meanrev]{Evidence of mean reversion?}
\begin{center}
\makebox[\linewidth]{
\includegraphics[scale=0.65]{figures/24June2024/meanrev_rph.pdf} 
}
\end{center}
\hyperlink{robustness}{\beamerbutton{Back}} 
\end{frame}

\end{document}

